\chapter{Аудит силы и сопротивления сквозного канала}
\label{ch:v10-throughflow-affinity-impedance-audit}

\section{Минимальный ориентированный цикл}

Стационарный сквозной ток нельзя реализовать одной обратимой стрелкой:
в двух состояниях стационарный ток сокращается условием подробного баланса.
Минимальный носитель ненулевой циркуляции содержит цикл. Рассмотрим три
состояния с прямой скоростью $2\kappa$ и обратной скоростью $\kappa$.
Генератор на вероятностях равен
\begin{equation}
 L_\circlearrowright=\kappa
 \begin{pmatrix}
 -3&1&2\\2&-3&1\\1&2&-3
 \end{pmatrix}.
 \label{eq:v10-throughflow-cycle-generator}
\end{equation}
Он имеет ранг $2$, одномерное ядро и стационарное состояние
\begin{equation}
 p_*=\frac13(1,1,1)^T.
\end{equation}
На каждом ребре сохраняется ненулевой ток
\begin{equation}
 J_{\rm edge}=\frac{2\kappa}{3}-\frac{\kappa}{3}=\frac{\kappa}{3}.
 \label{eq:v10-throughflow-cycle-current}
\end{equation}

\section{Сила цикла и производство энтропии}

Отношение прямой и обратной скоростей задаёт силу полного обхода
\begin{equation}
 F_\circlearrowright=3\log 2.
 \label{eq:v10-throughflow-cycle-affinity}
\end{equation}
Полное производство энтропии равно
\begin{equation}
 \sigma_\circlearrowright
 =3J_{\rm edge}\log2=\kappa\log2>0.
 \label{eq:v10-throughflow-cycle-entropy}
\end{equation}
Тем самым цикл реализует нужный сквозной режим: распределение стационарно,
но ориентированный ток и производство энтропии ненулевые.

Однако преобразование
\begin{equation}
 \kappa\mapsto c\kappa,
 \qquad t\mapsto\frac{t}{c}
 \label{eq:v10-throughflow-rate-clock-orbit}
\end{equation}
не меняет стационарное состояние и силу $3\log2$, тогда как ток и
производство энтропии умножаются на $c$. Карта отношения скоростей и
безразмерной эволюции имеет ранг $2$, ядро размерности $1$.

\section{Аудит кандидатов}

Проверены восемь источников силы и сопротивления: отношение скоростей КМС,
ориентация $K_{43}$, нормированная мера рождения, резонанс часов,
спектральная щель, хопфово число обвитий, резервуар физической материи и
космологический ток. Каждый оценён по шести условиям: происхождение силы,
абсолютная проводимость, циклический носитель, положительная энтропия,
внутренняя типизация и разрыв часовой орбиты.

Матрица имеет ранг $4$, полных проходов $0/8$. Лучший результат $4/6$
получает отношение скоростей КМС: оно фиксирует силу, но не общую величину
скоростей.

\begin{theorem}[Сила цикла не фиксирует скорость часов]
Ориентированный трёхвершинный цикл имеет единственное стационарное состояние,
ненулевой ток $\kappa/3$, силу $3\log2$ и положительное производство
энтропии $\kappa\log2$. Общая величина $\kappa$ остаётся свободной и задаёт
физический темп всей динамики.
\end{theorem}

\begin{nogocriterion}[Запрет подмены проводимости отношением скоростей]
КМС-отношение или топологическая ориентация могут выбрать направление и
безразмерную силу, но не абсолютную проводимость канала. Для физического
масштаба требуется внутренний механизм, фиксирующий сумму скоростей или
время релаксации независимо от внешних часов.
\end{nogocriterion}

\begin{protocol}[Статус силы и сопротивления]\begin{enumerate}
\item ориентированный стационарный цикл: \textbf{построен};
\item ненулевой ток и производство энтропии: \textbf{получены};
\item проверенных кандидатов: \textbf{$8$};
\item полных проходов: \textbf{$0/8$};
\item наибольший результат: \textbf{$4/6$};
\item абсолютная проводимость канала: \textbf{$0/1$};
\item физический масштаб часов: \textbf{$0/1$};
\item следующий гейт: \textbf{типизированное вложение хопфовского цикла в $K_{43}$}.
\end{enumerate}\end{protocol}

Машинный сертификат сохранён в
\path{s2t_v10_cell_birth_four_volume_throughflow_affinity_impedance_origin_audit_gate_results.json}.